

\subsection{Constraint-based Spatial Solving}
\label{subsec:p3_s7_1_4}
% - Continuous geometric/motion feasibility: identifying configurations or trajectories that satisfy geometric, dynamic, collision, and temporal constraints under a given goal.
% - CSPs, SAT/SMT, ILP/MILP, QP/NLP, IK solvers, feasible region modeling
% - Applications: manipulation planning, multi-robot coordination, assembly validation

Many robotics problems are not about finding \emph{a} plan but one that \emph{must} satisfy strict requirements: no collisions, no joint-limit violations, stable contacts, correct timing, sometimes guaranteed dynamics or safety margins. Constraint-based spatial solving states the problem as a collection of constraints, what must hold and what must not, and hands it to a solver that either returns a satisfying solution or proves that none exists. Compared with sampling- and heuristic-driven planning, this buys certificates of feasibility, an explanation of why a candidate fails, and, under an explicit model, optimality or bounded-suboptimality guarantees, which is what manipulation in clutter, multi-robot coordination in shared space, and assembly demand. Figure~\ref{fig:p3_s7_1_4_constraint_solving_abstract} gives the abstract formulation.

\begin{figure}[!ht]
    \centering
    \includegraphics[width=0.99\linewidth]{figures/part3/section7/714_constraint.pdf}
    \caption{Abstract task formulation of constraint-based spatial solving: (i) decision variables and constraints, (ii) a solver core that satisfies or optimizes them, and (iii) feasible trajectories, schedules, and assembly sequences, with feasibility or infeasibility certificates and optimality-gap bounds where supported by the solver and problem formulation. The output types are alternatives; failure of a local solver alone is not an infeasibility proof.}
    \label{fig:p3_s7_1_4_constraint_solving_abstract}
\end{figure}

\subsubsection{Task Definition}
Let $x=(\mathbf{q}_{1:T},\mathbf{u}_{1:T},\mathbf{z})$ collect continuous states and controls together with discrete mode and logic variables $\mathbf{z}$. Constraint-based solving seeks
\begin{equation}
\text{find } x\in\mathcal{X}\quad\text{s.t.}\quad g_i(x)\le 0,\ \ h_j(x)=0,\ \ \phi(\mathbf{z},x)=\top,
\end{equation}
where the inequalities $g_i$ encode geometric and dynamic feasibility (collision, joint, and dynamics limits), the equalities $h_j$ encode kinematics and contact conditions, and $\phi$ encodes task logic such as preconditions and effects. The optimization variant minimizes an objective $f(x)$ as well. The output is a feasible $x$ (a collision-free and dynamically admissible trajectory, a task-robot schedule, an assembly sequence), with an objective value or optimality gap when optimizing, or a certificate that no such $x$ exists.

\subsubsection{Methodology}
The solver families differ in the variables and constraints they handle and the guarantees they return; practical systems are rarely pure, combining discrete reasoning over actions and modes with continuous optimization over motions and contacts, and calling fast geometric oracles (collision checking, inverse kinematics) inside the solver loop.

\paragraph{Constraint Programming (CSP/CP).}
When the decisions are combinatorial (which object, which placement, which grasp) but each must pass a geometric test, constraint programming encodes the choices as variables under constraints and interleaves symbolic branching with propagation, pruning, and geometric checks. Lozano-P{\'e}rez and Kaelbling~\cite{LozanoPerez2014ConstraintTAMP} cast task-and-motion planning as constraint satisfaction over action skeletons, pruning symbolic branches as soon as a geometric feasibility test fails; Lagriffoul \etal~\cite{Lagriffoul2012ConstraintPropagation} propagate and filter continuous domains before any expensive collision or IK call, so that many failures become cheap inconsistencies; and Dantam \etal~\cite{Dantam2018IncrementalTAMP} interleave discrete search with incremental feasibility queries that reuse information across related constraints.

\paragraph{SAT/SMT Compilations.}
Compiling planning and feasibility into Boolean formulas, or into formulas over richer theories such as real arithmetic, puts clause learning, incremental solving, and conflict-driven pruning to work. dReal~\cite{Gao2013dReal} decides nonlinear real arithmetic up to a numerical tolerance ($\delta$-completeness), which gives certified answers where exact decidability is impossible; SMT-based planning for hybrid systems~\cite{Shoukry2016SMTPlanning} reasons first over an abstraction and adds constraints as counterexamples appear; SAT encodings of multi-agent pathfinding~\cite{Surynek2016SATMAPF} obtain optimal schedules on time-expanded graphs (the MAPF problem itself is treated in §\ref{subsec:p3_s7_1_3}); and SMT over linear real arithmetic extends the encodings to continuous time and to robustness against execution delays~\cite{Kolarik2023MAPFSMT,Tan2024RobustMAPF}.

\paragraph{Mixed-Integer Formulations (MICP/MILP).}
When discrete mode choices (which region, which contact mode, which precondition) interact with continuous geometry and dynamics, binary variables encode the nonconvex choice and convex constraints capture feasibility within a mode. IRIS~\cite{Deits2014IRIS} inflates large convex obstacle-free regions that serve as certified safe sets, turning collision avoidance into membership constraints with explicit margins; Graphs of Convex Sets~\cite{Marcucci2024GCSShortest} compile motion planning through clutter into a shortest-path problem over convex nodes and edges that convex optimization, with mixed-integer selection where needed, solves with certificates; and later work constructs convex sets directly in configuration space and in space-time~\cite{Marcucci2023SciRobotics,Dai2024CIRIS}. In each case a globally nonconvex feasibility problem is replaced by a sequence of convex ones without giving up the guarantee.

\paragraph{Trajectory Optimization.}
Even with the task fixed, an executable motion must satisfy kinematic limits, collision avoidance, smoothness, and often dynamics, which is a continuous, usually nonconvex, constrained optimization whose local optimum comes with optimality conditions and sensitivities. CHOMP~\cite{Zucker2013CHOMP} deforms an initial path by functional gradient descent on a signed-distance field; STOMP~\cite{Kalakrishnan2011STOMP} replaces the gradient with stochastic perturbations to handle non-smooth costs and local minima; TrajOpt~\cite{Schulman2014TrajOpt} solves a sequence of convex programs with hinge-loss collision penalties; and GPMP2~\cite{Mukadam2018GPMP2} writes planning as inference on a sparse factor graph with Gaussian-process priors, which gives scalability and incremental replanning. Increasingly these optimizers run inside the safe sets of IRIS and GCS, which supply initialization, constrain the search, or verify the result.

\paragraph{Inverse Kinematics as a Constraint Subproblem.}
``Is there a joint configuration that reaches this pose under limits and collisions?'' is the most frequent feasibility query in any hybrid planner, so IK solvers are built as fast oracles. IKFast~\cite{Diankov2010IKFast} derives closed-form solutions and generates optimized code for repeated real-time queries; TRAC-IK~\cite{Beeson2015TRACIK} combines complementary numerical strategies to raise the solve rate and avoid stalling the planner above it; and global IK by mixed-integer programming~\cite{Dai2019GlobalIK} returns certifiable solutions for full chains when a feasibility proof is needed.

\paragraph{Hybrid Logic--Geometric Planning.}
In long-horizon problems the discrete decision (what to do next) and the continuous one (how to do it) cannot be separated, and task-and-motion planners close the loop between them. Logic-Geometric Programming~\cite{toussaint2015lgp,Toussaint2017LGPcoop} couples symbolic action sequences to bound-constrained nonlinear programs through a multi-bound tree search, so that continuous feasibility and cost decide which symbolic branches are worth expanding; PDDLStream~\cite{garrett2020pddlstream} integrates black-box geometric generators (grasps, placements, motions) into symbolic planning by optimistic refinement, requesting continuous witnesses only when needed while remaining complete under suitable stream assumptions; R-LGP~\cite{Ly2023RLGP} biases the search with reachability guidance; and COAST~\cite{Vu2024COAST} uses constraints to communicate between task planner and stream generators, cutting wasted geometric proposals by an order of magnitude in planning time.

\subsubsection{Benchmarks}
Evaluation asks whether a method returns solutions that satisfy strict feasibility conditions and, increasingly, whether it can show that it did (margins, certificates, optimality gaps) under a comparable oracle budget of collision checks and IK calls. OMPL's benchmark suite and Planner Arena~\cite{Sucan2012OMPL} score planners on standardized problems with feasibility diagnostics and solution cost; the platform-independent TAMP benchmarks of Lagriffoul \etal~\cite{Lagriffoul2018TAMPBench} specify symbolic operators together with IK and collision feasibility so that a jointly feasible task-motion plan is required; and the assembly benchmarks RAMP~\cite{Collins2023RAMP} (real-world assemblies requiring geometric fit, stability under gravity, and grasp or fixture feasibility), ASAP~\cite{Tian2024ASAP} (large-scale 3D assembly sequences validated by physics-based feasibility checks), and FurnitureBench~\cite{Heo2023FurnitureBench} (long-horizon real-world furniture assembly scored by success rate and completion time) stress feasibility across geometry, stability, and sequencing at once. Multi-robot coordination under space-time collision constraints uses the MAPF benchmarks of §\ref{subsec:p3_s7_1_3}.

\subsubsection{Challenges and Trends}
The central difficulty is integrating a discrete task structure with continuous, highly nonconvex feasibility at scale and with certificates. Task-level choices (modes, grasps, placements, contact sets, schedules) interact combinatorially with geometric and dynamic constraints and fragment the feasible set; mixed-integer formulations grow expensive with the number of binary decisions~\cite{Deits2014IRIS,Marcucci2023SciRobotics}, trajectory optimizers return locally feasible motions that may violate a symbolic assumption such as the intended contact mode~\cite{Zucker2013CHOMP,Schulman2014TrajOpt}, and symbolic planners propose sequences that are geometrically unrealizable. Information flows poorly in both directions: continuous solvers rarely return explanations (active sets, separating hyperplanes, unsatisfiable cores) that the symbolic layer could use to prune, and symbolic abstractions are often too coarse to keep the continuous subproblems tractable~\cite{garrett2020pddlstream,toussaint2015lgp}. On the hybrid and assembly benchmarks this coupling, not raw search speed, is what separates planners: those that interleave symbolic branching with continuous oracles waste fewer feasibility checks and degrade more gracefully across domains.

Assurance under uncertainty is the second difficulty. Real systems need machine-checkable feasibility and fast infeasibility proofs, but nonlinear dynamics, contact, and hybrid interactions make strong certificates hard; convex surrogates and bounded-horizon relaxations buy speed at the price of completeness or tightness; execution requires real-time replanning with expensive oracles in the loop; and robustness to noise and model error calls for explicit risk and temporal-logic constraints~\cite{Blackmore2011ChanceConstraints,Raman2014STL,Ames2019CBF}. Evaluations still under-report certificate quality, oracle counts, and failure explanations.

The trend is toward structure with certificates: graphs of convex or geodesically convex safe sets that reduce global nonconvex feasibility to sequences of convex subproblems with explicit proofs~\cite{Marcucci2024GCSShortest,Dai2024CIRIS}, coupled to trajectory-optimization and MPC back-ends that return sensitivities and constraint activations for the search above them; chance constraints, temporal logic, and control barrier functions for persistent safety~\cite{Blackmore2011ChanceConstraints,Ames2019CBF}; GPU acceleration for real-time replanning; and learning that guides branching, cutting, or warm starts without discarding the guarantee~\cite{Ly2023RLGP,Vu2024COAST}.
